\section{Classical continuous laws as modelling tools}

Classical continuous distributions should not be learned as isolated formulas.
Each distribution is a mathematical model for a certain kind of random
phenomenon. The formula for the density is important, but the modelling idea is
more important: what mechanism or structure does the distribution represent?

This section introduces several important continuous laws at an interpretive
level. The aim is not to exhaust all their properties, but to understand why such
models appear and what questions they help answer.

\subsection*{Uniform distribution}

A random variable has a uniform distribution on \([a,b]\) if probability is spread
evenly over the interval. Its density is
\[
f_X(x)=
\begin{cases}
\frac1{b-a}, & a\leq x\leq b,\\
0, & \text{otherwise}.
\end{cases}
\]
The total area is
\[
\int_a^b\frac1{b-a}\,dx=1.
\]
For any subinterval \([c,d]\subseteq[a,b]\),
\[
P(c\leq X\leq d)=\frac{d-c}{b-a}.
\]
Thus probability is proportional to length.

The uniform distribution is appropriate when the model assumes that no region of
the interval is favoured over another. This assumption should be justified by the
context; it should not be used simply because the interval is known.

\subsection*{Exponential distribution}

The exponential distribution is commonly used to model waiting times. Its density
with rate parameter \(\lambda>0\) is
\[
f_X(x)=
\begin{cases}
\lambda e^{-\lambda x}, & x\geq0,\\
0, & x<0.
\end{cases}
\]
The density is largest at \(0\) and decreases as \(x\) increases. This means that
short waiting times are more likely than long waiting times, although long waits
remain possible.

The CDF is
\[
F_X(x)=1-e^{-\lambda x},\qquad x\geq0.
\]
Therefore
\[
P(X>t)=1-F_X(t)=e^{-\lambda t}.
\]
The parameter \(\lambda\) controls the scale of the waiting time. Larger
\(\lambda\) means that the density decreases more quickly and waiting times tend
to be shorter.

The exponential distribution is closely connected with the idea of a constant
instantaneous rate. In later applications, it appears as a model for time until a
random event occurs under a stable rate assumption.

\subsection*{Normal distribution}

The normal distribution is one of the central continuous models. Its density has
the form
\[
f_X(x)=\frac1{\sigma\sqrt{2\pi}}
\exp\left(-\frac{(x-
\mu)^2}{2\sigma^2}\right),
\qquad x\in\mathbb{R},
\]
where \(\mu\in\mathbb{R}\) and \(\sigma>0\).

The parameter \(\mu\) is the center of the distribution, and \(\sigma\) controls
its spread. The density is symmetric around \(\mu\). Values close to \(\mu\) have
higher density than values far away from \(\mu\).

The normal distribution is used when many small sources of variation combine or
when a distribution is approximately symmetric and bell-shaped. Its deeper
importance comes from limit theorems, especially the central limit theorem. This
connection will be developed later; for now, the normal distribution should be
read as a smooth, symmetric model for continuous variation.

\begin{remarkbox}
The normal density has no elementary antiderivative. Probabilities for normal
random variables are usually computed using tables, software, or standardized
normal distribution functions.
\end{remarkbox}

\subsection*{Gamma distribution}

The gamma distribution is a flexible model on \([0,\infty)\). It is often used for
positive quantities such as waiting times, lifetimes, or accumulated amounts. One
common parametrization has density
\[
f_X(x)=\frac{\lambda^\alpha}{\Gamma(\alpha)}x^{\alpha-1}e^{-\lambda x},
\qquad x>0,
\]
where \(\alpha>0\) is a shape parameter and \(\lambda>0\) is a rate parameter.
Here \(\Gamma(\alpha)\) is the gamma function, which generalizes factorials.

When \(\alpha=1\), the gamma distribution becomes the exponential distribution.
For larger shape parameters, the density can rise from zero, reach a peak, and
then decay. This makes the gamma family more flexible than the exponential
family.

A useful interpretation is that certain gamma distributions model the waiting
time until several events have occurred, not just the first one. Thus the gamma
law extends the waiting-time idea behind the exponential distribution.

\subsection*{Beta distribution}

The beta distribution is a flexible family of distributions on \([0,1]\). Its
density has the form
\[
f_X(x)=\frac{1}{B(\alpha,\beta)}x^{\alpha-1}(1-x)^{\beta-1},
\qquad 0<x<1,
\]
where \(\alpha>0\) and \(\beta>0\). The constant \(B(\alpha,\beta)\) normalizes
the density so that the total area is \(1\).

The beta distribution is useful for modelling proportions, probabilities, and
other quantities constrained between \(0\) and \(1\). Depending on the parameters,
the density can be concentrated near \(0\), near \(1\), near the middle, or even
near both endpoints.

For example, the uniform distribution on \([0,1]\) is a special beta
distribution with \(\alpha=1\) and \(\beta=1\). In that case,
\[
f_X(x)=1,
\qquad 0<x<1.
\]

\subsection*{Choosing a model}

A distribution should be chosen because its assumptions match the phenomenon
being modelled. The following questions are useful:
\begin{itemize}
    \item What values are possible? Are they all real numbers, only positive
    numbers, or values between \(0\) and \(1\)?
    \item Is the distribution expected to be symmetric or skewed?
    \item Are extreme values plausible or very unlikely?
    \item Is the variable a measurement, a waiting time, a proportion, or an
    accumulated quantity?
    \item Which assumptions justify the chosen family?
\end{itemize}

\begin{center}
\begin{tabular}{lll}
\toprule
Distribution & Support & Typical modelling idea \\
\midrule
Uniform & bounded interval & no region favoured \\
Exponential & \([0,\infty)\) & waiting time to first event \\
Normal & \(\mathbb{R}\) & symmetric continuous variation \\
Gamma & \([0,\infty)\) & positive waiting or accumulated quantity \\
Beta & \([0,1]\) & proportion or probability-like quantity \\
\bottomrule
\end{tabular}
\end{center}

\begin{summarybox}
Classical continuous laws are modelling tools. Their densities are formulas, but
the formulas should be connected to support, shape, parameters, and the mechanism
or interpretation that makes the model appropriate.
\end{summarybox}
